\GalSourcePageBoundary{082}{L0081}{53}{53}
nous obtiendrons ainsi, en tout, $P^{2}$~lettres. Si le nombre total
des lettres n'est pas épuisé, on prendra un troisième indice, en
sorte que
\[
a_{m,n,0}, \quad a_{m,n,1}, \quad a_{m,n,2}, \quad a_{m,n,3},\ \dots, \quad a_{m,n,P-1}
\]
soit, en général, un système de lettres conjointes; et l'on parviendra
ainsi à cette conclusion, que $N = P^{\mu}$, $\mu$~étant un certain
nombre égal à celui des indices différents dont on aura besoin. La
forme générale des lettres sera
\[
\GalPrintedError{a_{k_{1}, k_{1}, k_{3}, \dots, k_{\mu}}}{W10-PE001},
\]
$k_{1}$,~$k_{2}$, $k_{3}$,~\dots, $k_{\mu}$ étant des indices qui peuvent prendre chacun
les $P$~valeurs $0$,~$1$, $2$, $3$,~\dots, $P - 1$.

On voit aussi, par la manière dont nous avons procédé, que,
dans le groupe~$H$, toutes les substitutions seront de la forme
\[
\bigl[a_{k_{1}, k_{2}, k_{3} \dots k_{\mu}}, \quad
      \GalPrintedError{a_{\phi(k_{1}), \psi(k)_{2}, \chi(k_{3}) \dots, \sigma(k_{\mu})}}{W10-PE002}
\bigr],
\]
puisque chaque indice correspond à un système de lettres conjointes.

Si $P$ n'est pas un nombre premier, on raisonnera sur le groupe
de permutations de l'un quelconque des systèmes de lettres conjointes,
comme sur le groupe~$G$, en remplaçant chaque indice
par un certain nombre de nouveaux indices, et l'on trouvera
$P = R^{\alpha}$, et ainsi de suite; d'où enfin $N = p^{\nu}$, $p$~étant un nombre
premier.


\begin{center}\bfseries Des équations primitives de degré $p^{2}$.\end{center}

Arrêtons-nous un moment pour traiter de suite les équations
primitives d'un degré~$p^{2}$, $p$~étant premier impair. (Le cas de
$p = 2$ a été examiné.) Si une équation du degré~$p^{2}$ est soluble
par radicaux, supposons-la d'abord telle, qu'elle devienne non
primitive par une extraction de radical.

Soit donc $G$ un groupe primitif de $p^{2}$~lettres qui se partage
en $n$~groupes non primitifs conjugués à~$H$.

Les lettres devront nécessairement, dans le groupe~$H$, se ranger
